\documentclass{article}
\usepackage[utf8]{inputenc}
\usepackage[margin=1in,left=1.5in,includefoot]{geometry}
\usepackage{booktabs}
\usepackage{float}
% Header & Footer Stuff

\usepackage{fancyhdr}
\usepackage{wasysym}
\usepackage{amsmath}
\usepackage{tikz}
\usepackage{tabularx}
\usepackage{amssymb}
\usepackage{graphicx}
\pagestyle{fancy}
\lhead{MATH 220/199}
\rhead{March 24, 2026}
\fancyfoot{}
\lfoot{David Gallardo}
\fancyfoot[R]{}
%027601193
% The Main Document
\begin{document}

\begin{center}
 \LARGE\bfseries MATH 220/199 Week 9 Day 1
 \line(1,0){430}
\end{center}
\section*{Reminders}
\begin{enumerate}
    \item Office hours are held in Davenport 330 at 3pm on Thursdays.
    \item Do the Merit Quiz on the Merit Canvas page.
\end{enumerate}
\section*{The Assumptions of the MVT}
First, list the 2 assumptions/criteria which must be met in order to apply the MVT.
\vfill
\noindent
For each of the functions and intervals below, state whether all the assumptions of the MVT are met. If an assumption is not met, list it. If all assumptions are met, state the conclusion which you get from the MVT.
\begin{enumerate}
    \item $f(x) = |x-3|^2$
    \begin{enumerate}
        \item $[4, 6]$
        \vfill
        \item $[1, 4]$
        \vfill
    \end{enumerate}
    \item $g(x) = \sqrt{x^2-1}$
    \begin{enumerate}
        \item $[-2, 2]$
        \vfill
        \item $[3, 5]$
        \vfill
    \end{enumerate}
\end{enumerate}
\newpage
\section*{Exam Style Practice}
\begin{enumerate}
    \item Over the interval $[0, 3\pi]$, determine the coordinate(s) where $f(x) = 3\cos(\frac{x}{2})$ attains it's minimum value.
    \vfill
    \item Let $f(x) = \sec^2(x)$.
    \begin{enumerate}
        \item Over the open interval $(-\pi, \pi)$, determine the intervals where $f(x)$ is decreasing and increasing (you may present your answer in the form of a sign chart).
        \vfill
        \item Determine if $f(x)$ attains maximum/minimum values over $(-\pi, \pi)$. If it does, list them.
        \vfill
    \end{enumerate}
    \item Let $g(x) = \begin{cases}
        (x-1)^2+1 & x \ne 1\\
        5 & x=1
    \end{cases}$. Determine the minimum and maximum values attained by $g(x)$ over the interval $[0, 3]$ as well as where they occur.
    \vfill
\end{enumerate}

\end{document}
